
\begin{proposition}[The effective conductor divisor on a minimal resolution]\label{cdp:st:A.5.3}
Let $S$ be a compact reduced hypersurface in a smooth threefold, let $\eta:\coprod_i T_i\to S$ be its normalization, and let $\pi_i:Z_i\to T_i$ be the minimal resolutions. Put $h_i=\eta\pi_i$. There are canonical effective integral divisors $D_i$ on $Z_i$ with
\[
 \omega_{Z_i}(D_i)\simeq h_i^*\omega_S.
 \tag{A.25}\label{cdp:eq:A.25}
\]
The nonexceptional part of $D_i$ is exactly the divisorial conductor, with its scheme-theoretic multiplicities. When $S\subset Y$ and $H^2(Y,\Q)=0$, every $K_{Z_i}+D_i$ has rational first Chern class zero. A ruled $Z_i\to B$ with $g(B)\ge2$ then has
\[
 D_i=2C+V,
 \tag{A.26}\label{cdp:eq:A.26}
\]
where $C$ is a section and $V$ is vertical. In the case where $C$ is $\pi_i$-exceptional, $\eta^*(L|_S)|_{T_i}$ is trivial for every $L\in\Pic(Y)$.
\end{proposition}
\begin{proof}
Finite duality gives $\omega_{T_i}=\mathfrak c_i\otimes\eta^*\omega_S$ as rank-one reflexive sheaves, where $\mathfrak c_i$ is the conductor viewed as an ideal on $T_i$. Locally this is the equality
$\Hom_R(B,R)=\{q:qB\subset R\}=\mathfrak c$ for a reduced Gorenstein ring $R$ and its finite normalization $B$, tensored with the invertible dualizing module. The equality uses the meromorphic multiplication identification as in \ref{cdp:st:A.2.1}. The conductor is a maximal Cohen--Macaulay module because it is this dual of the maximal Cohen--Macaulay normalization over the Gorenstein ring. On a normal surface a rank-one maximal Cohen--Macaulay module is reflexive, so the displayed equality is determined by its codimension-one orders. Those orders define an effective conductor divisor $D_i^0$ on $T_i$.

Off the exceptional set the canonical identification pulls back to $\omega_{Z_i}(D_i^{\rm str})\simeq h_i^*\omega_S$, where $D_i^{\rm str}$ denotes strict transform. The meromorphic extension of this line-bundle identification gives integral exceptional coefficients $a_j$ with
$K_{Z_i}+D_i^{\rm str}+\sum_j a_jE_j=h_i^*K_S$.
We verify their signs. For each exceptional irreducible curve $E_j$, minimality excludes a smooth rational $(-1)$-curve. Adjunction therefore gives
$K_{Z_i}\cdot E_j=2p_a(E_j)-2-E_j^2\ge0$;
also $D_i^{\rm str}\cdot E_j\ge0$. The exceptional intersection matrix $M=(E_j\cdot E_k)$ is negative definite with nonnegative off-diagonal entries, and
\[
 M a=-v,\qquad v_j=(K_{Z_i}+D_i^{\rm str})\cdot E_j\ge0.
\]
To see that $a\ge0$ without an unproved sign convention, write $a=a_+-a_-$ with disjoint nonnegative coordinate supports. If $a_-\ne0$, then
$a_-^tMa=a_-^tMa_+-a_-^tMa_->0$,
whereas $Ma=-v$ gives $a_-^tMa\le0$. This contradiction proves $a_j\ge0$. Thus $D_i=D_i^{\rm str}+\sum_j a_jE_j$ proves \eqref{cdp:eq:A.25}. The conductor duality and effective exceptional comparison also appear, for projective surfaces with trivial dualizing sheaf, in \cite[Lemma~1.1 and Corollary~1.4]{Furushima1993}; the argument here supplies the maps and does not use that paper's additional K3-type hypothesis $h^1(\OO_S)=0$.

For a divisor in $Y$, adjunction gives $\omega_S=(K_Y\otimes\OO_Y(S))|_S$, proving the numerical assertion. On a general ruling fibre $f$, $K_{Z_i}\cdot f=-2$, so $D_i\cdot f=2$. Every horizontal irreducible curve has normalization of genus at least $g(B)\ge2$. Such a curve $C_0$ cannot occur in $D_i$ with coefficient one: otherwise
\[
 0=(K_{Z_i}+D_i)\cdot C_0
   =2p_a(C_0)-2+(D_i-C_0)\cdot C_0>0.
\]
Every horizontal coefficient is consequently at least two. The degree-two intersection with $f$ forces exactly one horizontal component, coefficient two and degree one over $B$, giving \eqref{cdp:eq:A.26}. The degree-one finite morphism $C\to B$ is an isomorphism because $B$ is normal; thus $C$ is a smooth section.

Finally, integral $H^2$ of a smooth ruled surface and of its blowups is torsion-free, and its topologically trivial line bundles are precisely the pullbacks of $\Pic^0(B)$. These facts follow from the projective-bundle and blowup cohomology formulas and from $p_*\OO_{Z_i}=\OO_B$, $R^1p_*\OO_{Z_i}=0$. Restriction to the section $C$ is inverse to this Picard pullback. The line $h_i^*L$ has torsion, hence zero, integral first Chern class. When $C$ is exceptional, its restriction to $C$ is trivial because $h_i$ maps $C$ to a point. Therefore $h_i^*L$ is trivial on $Z_i$. Normality gives $(\pi_i)_*\OO_{Z_i}=\OO_{T_i}$, and the projection formula descends this triviality to $T_i$. This proves the last assertion.
\end{proof}


